\section{Ferromagnetic laws and their clocks}\label{ch:fmlaws}



\subsection{Definition and background}\label{fm:sec:def}

\begin{definition}\label{fm:def:ferro}
A real random variable $X$ is \emph{ferromagnetic} if it is symmetric, $\E e^{b_0X^2}<\infty$ for some $b_0>0$, and its moment generating function $L_X(z)=\E e^{zX}$, which is then entire, vanishes only on the imaginary axis. Equivalently, every zero of the characteristic function $\varphi_X(\theta)=\E e^{i\theta X}=L_X(i\theta)$ is real.
\end{definition}

These are Newman's laws of type $\mathcal L$ \cite{Newman75}.

The name comes from statistical mechanics. Lee and Yang proved that the partition function of an Ising ferromagnet, with spins $\pm1$ and a common complex field, vanishes only when the field is purely imaginary \cite{LeeYang}, and Newman showed that a symmetric single spin law with Gaussian moments of all orders has the same property in every ferromagnetic system exactly when its own moment generating function has only imaginary zeros \cite{Newman74}; Lieb and Sokal gave another proof and extended the theorem to spins with several components \cite{LiebSokal81}. Section~\ref{rz:sec:polya} recalls these results, and Proposition~\ref{rz:prop:newmanthorin} restates Newman's theorem in the language of Thorin measures. Two smaller classes of single spin laws carry the correlation inequalities of the Ising model. The Griffiths Simon class consists of the weak limits of suitably scaled sums of ferromagnetically coupled Ising spins; it contains the laws with density proportional to $e^{-ax^4-bx^2}$, $a>0$, and its laws inherit the Lee and Yang property and the inequalities of Griffiths, Kelly and Sherman and of Griffiths, Hurst and Sherman \cite{SimonGriffiths73,GHS70}. The class of Ellis, Monroe and Newman consists of the laws with density $e^{-V(x)}$, $V$ even and $V'$ convex on $[0,\infty)$, and their weak limits; the Griffiths, Hurst and Sherman inequality holds for ferromagnets with such single spin laws \cite{EMN76}.

\subsection{The Hadamard form and the cumulants}\label{fm:sec:hadamard}

\begin{theorem}[Newman \cite{Newman75}]\label{fm:thm:hadamard}
Let $X$ be symmetric with $\E e^{b_0X^2}<\infty$ for some $b_0>0$. Then $X$ is ferromagnetic if and only if
\begin{equation}\label{fm:eq:had}
L_X(s)=e^{bs^2}\prod_k\Big(1+\frac{s^2}{c_k^2}\Big),\qquad b\ge0,\quad c_k>0,\quad\sum_kc_k^{-2}<\infty,
\end{equation}
with each $c_k$ repeated according to the multiplicity of the zero $ic_k$. If $\E e^{bX^2}<\infty$ for every $b>0$, then $b=0$.
\end{theorem}

\begin{proof}
Since $|L_X(z)|\le\E e^{|\re z||X|}\le Ce^{|z|^2/4b_0}$, the function $\varphi_X$ is entire of order at most $2$; it is even and real on $\R$, and $\varphi_X(0)=1$. Suppose that its zeros are real, at $\pm c_k$ with $c_k>0$. Hadamard's theorem, with the primary factors of genus $2$ paired over $\pm c_k$, gives $\varphi_X(\theta)=e^{-b'\theta^2}\prod_k(1-\theta^2/c_k^2)e^{\theta^2/c_k^2}$ with $b'$ real and $\sum_kc_k^{-4}<\infty$, so that, for real $s$,
\[
\frac{\log L_X(s)}{s^2}=b'+\sum_k\frac1{c_k^2}\,\frac{\log(1+x_k)-x_k}{x_k},\qquad x_k=\frac{s^2}{c_k^2}.
\]
The function $(\log(1+x)-x)/x$ decreases from $0$ to $-1$ on $(0,\infty)$, so by monotone convergence the sum tends to $-\sum_kc_k^{-2}$ as $s\to\infty$. If this sum diverged, $L_X(s)$ would tend to $0$, which is impossible because $L_X(s)=\E\cosh(sX)\ge1$. Hence $\sum_kc_k^{-2}<\infty$, the factors $e^{\theta^2/c_k^2}$ can be collected into the exponential, and \eqref{fm:eq:had} holds with $b=b'-\sum_kc_k^{-2}$. Then $\log L_X(s)=bs^2+o(s^2)$, since $\log(1+x)/x$ is bounded by $1$ and tends to $0$, and $L_X\ge1$ forces $b\ge0$. The converse is clear, since the right side of \eqref{fm:eq:had} vanishes only at $\pm ic_k$. If $\E e^{bX^2}<\infty$ for every $b$, then $L_X(s)\le C_\epsilon e^{\epsilon s^2}$ for every $\epsilon>0$, so $b=0$, as in the proof of Proposition~\ref{rz:prop:newmanthorin}.
\end{proof}



\begin{corollary}[Newman \cite{Newman75}]\label{fm:cor:cum}
If $X$ is ferromagnetic, with \eqref{fm:eq:had}, then $\operatorname{Var}X=2b+2\sum_kc_k^{-2}$, and for $m\ge2$
\begin{equation}\label{fm:eq:cum}
\frac{\kappa_{2m}(X)}{(2m)!}=\frac{(-1)^{m+1}}m\sum_kc_k^{-2m}.
\end{equation}
So the even cumulants alternate in sign: $\kappa_4\le0$, $\kappa_6\ge0$, and so on. Moreover $\E e^{sX}\le e^{s^2\operatorname{Var}X/2}$ for real $s$, and $\E X^4\le3(\operatorname{Var}X)^2$.
\end{corollary}

\begin{proof}
The first two statements are the expansion of $\log L_X(s)=bs^2+\sum_k\log(1+s^2/c_k^2)$. The bound on $\E e^{sX}$ follows from $\log(1+x)\le x$, and the bound on $\E X^4=\kappa_4+3\kappa_2^2$ from $\kappa_4\le0$.
\end{proof}

For the P\'olya law the alternation holds with no hypothesis: by the formula for $\log G$ in Section~\ref{ch:heat}, $(-1)^{n-1}\kappa_{2n}(X)=\frac{(2n)!}{n!}\int_0^\infty t^{n-1}W(t)\,dt>0$ for every $n$ (Section~\ref{ch:ggc}), and under RH this is \eqref{fm:eq:cum} with $c_k=\gamma_k$ and $b=0$.

\subsection{The clock of a ferromagnetic law}\label{fm:sec:clock}



\begin{theorem}\label{fm:thm:dual}
Let $X$ be ferromagnetic, with \eqref{fm:eq:had}, and put
\[
T_X=b+\sum_k\frac{e_k}{c_k^2},\qquad Y_X=Z\sqrt{2T_X},
\]
the series converging almost surely and in mean.
\begin{enumerate}
\item[(i)] $\E e^{-uT_X}=1/L_X(\sqrt u)$ for $u\ge0$, and $T_X$ is a GGC with drift $b$ and Thorin measure $\sum_k\delta_{c_k^2}$, one unit atom for each pair of zeros $\pm ic_k$.
\item[(ii)] $\E e^{i\theta Y_X}=1/L_X(\theta)$ for real $\theta$, so that $(\varphi_X,\varphi_{Y_X})$ is a van Dantzig pair, and $Y_X\overset d=\sqrt{2b}\,Z+\sum_kL_k/c_k$.
\item[(iii)] The moment generating function of $Y_X$ is $1/\varphi_X(t)$ for $|t|<c_1$ and infinite for $|t|\ge c_1$, unless $X$ is Gaussian.
\item[(iv)] Unless $X$ is degenerate, $Y_X$ has a P\'olya frequency density of infinite order.
\item[(v)] $\kappa_m(T_X)=(-1)^{m+1}m!\,\kappa_{2m}(X)/(2m)!$ for $m\ge1$; in particular $\E T_X=\frac12\operatorname{Var}X$ and $\operatorname{Var}T_X=-\kappa_4(X)/12$.
\end{enumerate}
\end{theorem}

\begin{proof}
(i) $\E e^{-ue_k/c_k^2}=(1+u/c_k^2)^{-1}$; multiply over $k$ and compare with \eqref{fm:eq:had} at $s=\sqrt u$. (ii) Conditionally on $T_X$, $Y_X$ is centred normal with variance $2T_X$, so $\E e^{i\theta Y_X}=\E e^{-\theta^2T_X}=1/L_X(\theta)$, which is also the characteristic function of $\sqrt{2b}\,Z+\sum_kL_k/c_k$, since $\E e^{i\theta L}=(1+\theta^2)^{-1}$. (iii) $\E e^{tY_X}=\E e^{t^2T_X}=e^{bt^2}\prod_k(1-t^2/c_k^2)^{-1}$ while $t^2<c_1^2$, and $\E e^{t^2e_1/c_1^2}=\infty$ for $t^2\ge c_1^2$. (iv) The Laplace density $\frac12e^{-|x|}$ and the normal densities are P\'olya frequency functions of infinite order, and the class is closed under convolution and under the limits involved here. Equivalently, the bilateral Laplace transform of $Y_X$ is the reciprocal of $\varphi_X(s)=e^{-bs^2}\prod_k(1-s^2/c_k^2)$, which has Schoenberg's form \cite{Schoenberg51}. (v) For $u$ near $0$, $\log\E e^{-uT_X}=-\log L_X(\sqrt u)=-\sum_m\kappa_{2m}(X)u^m/(2m)!$, and comparing coefficients with $\sum_m\kappa_m(T_X)(-u)^m/m!$ gives the formula.
\end{proof}

Part~(i) is Proposition~\ref{rz:prop:newmanthorin} with the Gaussian factor allowed. For the P\'olya law under RH, $b=0$, $c_k=\gamma_k$ and $T_X$ is the clock $T$ of Proposition~\ref{prop:rigid}, and $(\varphi_X,\varphi_{Y_X})$ is the van Dantzig pair of Section~\ref{ch:clock}, with $Y_X=Z\sqrt{2T}$ the variable $Y_{1/2}$ of Proposition~\ref{sm:prop:cf}. Without RH the clock $T$ still exists (Theorem~\ref{thm:ch12:idsd}), and so does the pair; what fails without RH is the product form of the clock. Konstantopoulos, Patie and Sarkar study van Dantzig pairs of this kind through the Laguerre P\'olya class and P\'olya frequency functions, and connect them with the Lee and Yang property \cite{KPS24}. The correspondence runs one way only: every ferromagnetic law gives a P\'olya frequency law, but $\varphi(s)=1-s^2$, the reciprocal of the bilateral Laplace transform of the Laplace law, is not a characteristic function.

Part~(v) holds for the P\'olya law with no hypothesis, with the clock $T$ of Theorem~\ref{thm:ch12:idsd} in place of $T_X$: the identity $\log\E e^{-uT}=-\log G(u)=-\log\E e^{\sqrt uX}$ is all the proof uses. With the cumulants of Table~\ref{fm:tab:cum} it gives $\E T=0.0231049931$ and $\operatorname{Var}T=-\kappa_4/12=3.71725993\times10^{-5}$, the constants of Section~\ref{ch:clock}.

\subsection{Examples}\label{fm:sec:examples}

Table~\ref{fm:tab:examples} lists the basic ferromagnetic laws with their clocks. The variables are those of Biane, Pitman and Yor \cite{BPY01}: $C$ is the exit time of Brownian motion from $(-1,1)$, with $\E e^{-uC}=1/\cosh\sqrt{2u}$, and $J$ is the hitting time of $1$ by a Bessel process of dimension $3$ started at $0$, with $\E e^{-uJ}=\sqrt{2u}/\sinh\sqrt{2u}$, as in Section~\ref{bp:sec:mellin}. The first two rows are the classical van Dantzig pairs of Section~\ref{ch:clock}.

\begin{table}[htbp]
\centering\small\setlength{\tabcolsep}{4.5pt}
\begin{tabular}{@{}lllll@{}}
\toprule
law of $X$ & $L_X(s)$ & $c_k$ & clock $T_X$ & $Y_X$\\
\midrule
fair sign $\varepsilon$ & $\cosh s$ & $(k-\frac12)\pi$ & $C/2$ & hyperbolic secant\\
uniform on $[-1,1]$ & $\sinh s/s$ & $k\pi$ & $J/2$ & logistic, scale $1/\pi$\\
$c_\nu(1-x^2)^{\nu-1/2}$, $\nu>-\frac12$ & $\Gamma(\nu+1)(2/s)^\nu I_\nu(s)$ & $j_{\nu,k}$ & $\frac12\tau_\nu$ & Bessel type\\
$N(0,\sigma^2)$ & $e^{\sigma^2s^2/2}$ & none & $\sigma^2/2$ & $N(0,\sigma^2)$\\
$e^{-2\pi\cosh2x}/K_0(2\pi)$ & $K_{s/2}(2\pi)/K_0(2\pi)$ & $2\nu_k$ & $T_P/4$ & $Z\sqrt{T_P/2}$\\
P\'olya, under RH & $\xi(\frac12+s)/\xi(\frac12)$ & $\gamma_k$ & $T$ & $Z\sqrt{2T}$\\
\bottomrule
\end{tabular}
\caption{Ferromagnetic laws and their clocks. Here $j_{\nu,k}$ are the positive zeros of the Bessel function $J_\nu$, and $\tau_\nu$ is the hitting time of $1$ by a Bessel process of dimension $2\nu+2$ started at $0$. In the fifth row $\nu_k$ are the zeros of $\nu\mapsto K_{i\nu}(2\pi)$ and $T_P$ is the clock of Theorem~\ref{thm:biane:indexP}.}
\label{fm:tab:examples}
\end{table}

The third row rests on Kent's theorem for Bessel processes \cite{Kent78}: for the Bessel process of dimension $2\nu+2$ started at $0$, the hitting time of $1$ satisfies $\E e^{-u\tau_\nu}=(\sqrt{2u})^\nu/(2^\nu\Gamma(\nu+1)I_\nu(\sqrt{2u}))$, so that $\E e^{-u\tau_\nu/2}=1/L_X(\sqrt u)$ and $T_X=\frac12\tau_\nu=\sum_ke_k/j_{\nu,k}^2$. The zeros $j_{\nu,k}$ are real for $\nu>-1$, so every symmetric beta law $c_\nu(1-x^2)^{\nu-1/2}$ with $\nu>-\frac12$ is ferromagnetic. The family contains the arcsine law ($\nu=0$), the uniform law ($\nu=\frac12$) and the semicircle law ($\nu=1$), and it tends to the law of $\varepsilon$ as $\nu\downarrow-\frac12$, where $\tau_{-1/2}=C$.

\subsection{Closure}\label{fm:sec:closure}

Products of functions of the form \eqref{fm:eq:had} have the same form, so sums of independent ferromagnetic variables are ferromagnetic. The class is closed under convergence in law, and this is what the Gaussian factor in Definition~\ref{fm:def:ferro} is for: the laws of $n^{-1/2}\sum_{k\le n}\varepsilon_k$ have Gaussian moments of all orders, and their limit $N(0,1)$ does not.

\begin{theorem}\label{fm:thm:closed}
Let $X_n$ be ferromagnetic and suppose that $X_n\to X$ in law. Then $X$ is ferromagnetic, $\operatorname{Var}X_n\to\operatorname{Var}X$, and $L_{X_n}\to L_X$ locally uniformly on $\C$.
\end{theorem}

\begin{proof}
Let $\sigma_n^2=\operatorname{Var}X_n$. By Corollary~\ref{fm:cor:cum}, $\E X_n^4\le3\sigma_n^4$, so the Paley Zygmund inequality applied to $X_n^2$ gives $P(X_n^2>\sigma_n^2/4)\ge(\frac34)^2\sigma_n^4/\E X_n^4\ge\frac3{16}$. A sequence that converges in law is tight, so the $\sigma_n$ are bounded, say by $\sigma$. By Corollary~\ref{fm:cor:cum} again, $|L_{X_n}(z)|\le L_{X_n}(\re z)\le e^{\sigma^2(\re z)^2/2}$, so the $L_{X_n}$ form a locally bounded, hence normal, family of entire functions, and for each real $s$ the variables $e^{sX_n}$ are uniformly integrable, since their second moments are bounded. Hence $\E e^{sX_n}\to\E e^{sX}$ for real $s$, which gives $\E e^{sX}\le e^{\sigma^2s^2/2}$ and, averaging over $s=\sqrt{2b}Z$, $\E e^{bX^2}<\infty$ for $b<1/2\sigma^2$. So $L_X$ is entire, every locally uniform limit point of the $L_{X_n}$ agrees with it on $\R$, and by Vitali's theorem $L_{X_n}\to L_X$ locally uniformly. Each $L_{X_n}$ has no zero in the half plane $\re z>0$, and $L_X(0)=1$, so by Hurwitz's theorem $L_X$ has no zero there either; by symmetry it has none in $\re z<0$. The fourth moments are bounded, so the $X_n^2$ are uniformly integrable and the variances converge.
\end{proof}



\begin{proposition}\label{fm:prop:tilt}
Let $X$ be ferromagnetic with law $\rho$, and let $\E e^{b_0X^2}<\infty$.
\begin{enumerate}
\item[(i)] For every $a>0$ the law $\cosh(ax)\,\rho(dx)/\E\cosh(aX)$ is ferromagnetic.
\item[(ii)] For every $0<\lambda<b_0$ the law $e^{\lambda x^2}\rho(dx)/\E e^{\lambda X^2}$ is ferromagnetic.
\end{enumerate}
\end{proposition}

\begin{proof}
(i) The moment generating function of the tilted law is proportional to $L_X(z+a)+L_X(z-a)$. For $w$ with $\re w>0$ and $a>0$ we have $|w+a|>|w-a|$. Applied to $w=z\pm ic_k$ this gives $|c_k^2+(z+a)^2|>|c_k^2+(z-a)^2|$ when $\re z>0$, and $|e^{b(z+a)^2}|=e^{4ab\re z}|e^{b(z-a)^2}|$. So for $\re z>0$ every factor of \eqref{fm:eq:had} has larger modulus at $z+a$ than at $z-a$, and $L_X(z+a)+L_X(z-a)\ne0$ unless $L_X$ is constant, in which case the sum is $2$. The sum is even in $z$, so it has no zeros with $\re z<0$ either, and the tilted law satisfies the moment condition because $\cosh(ax)\le e^{a|x|}$. (ii) By (i), applied $n$ times, the law with density proportional to $\cosh(x\sqrt{2\lambda/n})^n$ with respect to $\rho$ is ferromagnetic. Since $\cosh y\le e^{y^2/2}$, these densities are dominated by $e^{\lambda x^2}$, which is $\rho$ integrable, and they converge to $e^{\lambda x^2}$, so the laws converge in total variation to $e^{\lambda x^2}\rho(dx)/\E e^{\lambda X^2}$. Theorem~\ref{fm:thm:closed} applies.
\end{proof}

Part~(i) is the one spin case of the Lee and Yang theorem, and for the P\'olya law, whose transform has genus $1$, it is P\'olya's Hilfssatz~II of Section~\ref{rz:sec:polya}. For the P\'olya law, part~(ii) is the forward half of de Bruijn's theorem \cite{deBruijn}: multiplying $p$ by $e^{\lambda x^2/4}$ is the de Bruijn Newman flow at time $\lambda$ (Section~\ref{ch:dbn}).


\begin{corollary}\label{fm:cor:limit}
RH holds if and only if the P\'olya law is a limit in law of ferromagnetic laws.
\end{corollary}

\begin{proof}
If RH holds, the P\'olya law is ferromagnetic by \eqref{eq:ch2:mgf}, and it is its own limit. Conversely, a limit in law of ferromagnetic laws is ferromagnetic by Theorem~\ref{fm:thm:closed}, and then every zero of $\Xi$ is real.
\end{proof}

Let $\delta$ be the infimum of the total variation distances $\sup_A|P(X\in A)-P(Y\in A)|$ from the P\'olya variable $X$ to ferromagnetic variables $Y$; when $Y$ has a density $f$ the distance is $\operatorname{TV}(f,p)=\frac12\int|f-p|$. If RH holds, $\delta=0$; conversely, ferromagnetic laws that converge to the P\'olya law in total variation converge in law, so $\delta=0$ implies RH.
